\documentclass[11pt,a4paper]{article}
\usepackage[T1]{fontenc}
\usepackage[utf8]{inputenc}
\usepackage{lmodern}
\usepackage{amsmath,amssymb,amsthm}
\usepackage[margin=1.15in]{geometry}
\usepackage{booktabs}
\usepackage{enumitem}
\usepackage[hidelinks]{hyperref}
\emergencystretch=3em

\newtheorem{theorem}{Theorem}[section]
\newtheorem{proposition}[theorem]{Proposition}
\newtheorem{lemma}[theorem]{Lemma}
\newtheorem{corollary}[theorem]{Corollary}
\newtheorem*{theoremA}{Theorem A}
\newtheorem*{theoremB}{Theorem B}
\newtheorem*{corollaryC}{Corollary C}
\theoremstyle{definition}
\newtheorem{definition}[theorem]{Definition}
\newtheorem{remark}[theorem]{Remark}
\newtheorem{example}[theorem]{Example}

\newcommand{\Cat}{\mathsf{Cat}}
\newcommand{\sSet}{\mathsf{sSet}}
\newcommand{\Fl}{\operatorname{Fl}}
\newcommand{\Loc}[1]{\mathrm{Loc}\,(#1)}
\newcommand{\Winf}{W_{\infty}}
\newcommand{\Wpi}{W_{\pi}}
\newcommand{\Wcirc}{W_{\mathrm{circ}}}
\newcommand{\Wtor}{W_{\mathrm{tor}}}
\newcommand{\Wfin}{W_{\mathrm{fin}}}
\newcommand{\Wctbl}{W_{\mathrm{ctbl}}}
\newcommand{\Tor}{\operatorname{Tor}}
\newcommand{\Z}{\mathbb{Z}}
\newcommand{\Q}{\mathbb{Q}}
\newcommand{\N}{\mathbb{N}}
\newcommand{\id}{\mathrm{id}}
\newcommand{\colim}{\operatorname*{colim}}
\newcommand{\Ho}{\operatorname{Ho}}
\newcommand{\Pione}{\Pi_1}

\title{Weak fundamental localizers that are not fundamental localizers,
and Corollaire~7.7.2 of Grothendieck's \emph{Les D\'erivateurs}\thanks{AI-produced.
The four classes and the sufficiency in Theorem~B are due to runs of GPT-6 Astra (OpenAI,
ChatGPT Pro mode); the necessity, some complements and the text are by Claude Opus~5.5
(Anthropic) at effort max; Proposition~\ref{prop:ctbl} and two further complements are by
Claude Fable~5.2 (Anthropic) at effort max; see Section~\ref{sec:method} for who proved, wrote
and checked what. No human has checked this text. To the extent that any
rights exist in it, they are dedicated to the public domain under CC0~1.0.}}
\author{}
\date{25 November 2026}

\begin{document}
\maketitle

\begin{abstract}
In Chapter~XVI of \emph{Les D\'erivateurs} Grothendieck axiomatized the classes $W$ of weak
equivalences between small categories by a saturation axiom, the asphericity of categories
with a final object, and Quillen's Theorem~A in an absolute form, Loc~(3). He wrote that a
relative form, Loc~(3~bis), was the more important one, and in Corollaire~7.7.2 he derived
it from three further axioms: gluing along Mayer--Vietoris squares, Loc~($6'$a); stability
under filtered colimits, Loc~(8); and accessibility, Loc~(9). A footnote adds: ``Loc (3 bis)
pas prouv\'e, probablement faux''. Classes satisfying the absolute form are now called weak
fundamental localizers (Maltsiniotis), classes satisfying the relative form fundamental
localizers, and Cisinski proved that the weak homotopy equivalences form the smallest class
of either kind. We show that the absolute form does not imply the relative one: there are
weak fundamental localizers that are not fundamental localizers. The smallest of our witnesses is a
functor between groupoids with two objects over the group of order two; other examples are
strongly saturated and satisfy Loc~(8) and Loc~(9), or are accessible and satisfy
Loc~($6'$a). Conversely, a weak fundamental localizer is a fundamental localizer if and only
if it satisfies Loc~($6'$a) and Loc~(8); this direction rests on Cisinski's correspondence
between fundamental localizers and localizers of simplicial sets. So Corollaire~7.7.2 holds,
and holds without Loc~(9): the footnote's ``probablement faux'' is not borne out. Neither of
the other two hypotheses can be omitted, and the doubt becomes true as soon as one of them is
dropped; nor can Loc~(8) be weakened to stability under countable colimits. In all our examples the
aspherical categories are exactly the weakly contractible ones, so the difference between the
two axiom systems is invisible to asphericity. The four classes and the sufficiency of the two axioms
were found by GPT-6 Astra (OpenAI, ChatGPT Pro mode); their necessity and the text are by
Claude Opus~5.5 (Anthropic) at effort max, and the example showing that Loc~(8) cannot be
weakened is by Claude Fable~5.2 (Anthropic) at effort max. The proofs were checked by Claude
Opus~5.5 at efforts max and high and by Claude Fable~5.2 at effort max. No human has checked
them.
\end{abstract}

\section{Introduction}\label{sec:intro}

In Chapter~XVI of the manuscript \emph{Les D\'erivateurs} \cite{Der}, Grothendieck returns to
the axioms for a class $W$ of functors between small categories that is to play the r\^ole
of the weak homotopy equivalences. His Définition~1.1 calls \emph{localiseur fondamental} a
class satisfying three axioms: a saturation axiom Loc~(1), the asphericity of categories
with a final object Loc~(2), and Loc~(3), Quillen's Theorem~A in an absolute form: a functor
whose comma categories are all aspherical belongs to $W$. On the next page he states the
relative form Loc~(3~bis), of which Loc~(3) is the special case of a functor over its own
target, and in footnote~1 he proposes to exchange the two numberings in a final version,
because Loc~(3~bis) ``m'appara\^it comme le plus important des deux'' \cite[XVI, p.~1]{Der}.
At the end of the chapter, in Corollaire~7.7.2 (manuscript page~118), he derives Loc~(3~bis)
from three further axioms, and adds footnote~56: ``Loc (3 bis) pas prouv\'e, probablement
faux.'' The preceding footnote~55 names the difficulty: the identification of $W$ with the
localizer of the derivator $\mathrm{HOT}_W$, on which the argument rests, had been proved
under Loc~(3~bis) (``Gare au cercle vicieux'').

The terminology now in use follows footnote~1. Cisinski \cite[D\'ef.~3.3.2]{Cis06} calls
\emph{fundamental localizer} a class satisfying the relative form, and, following
Maltsiniotis \cite[1.1.2]{Mal05}, \emph{weak fundamental localizer} a class satisfying the
absolute form \cite[6.1.12]{Cis06}. He observes that a large part of Grothendieck's
homotopy theory, for instance the theory of test categories and test functors, only uses the
axioms of weak fundamental localizers, and he proves that the class $\Winf$ of functors whose
nerve is a weak homotopy equivalence is the smallest weak fundamental localizer
\cite[Th.~6.1.18]{Cis06}, as it is the smallest fundamental localizer
\cite[Cor.~4.2.19]{Cis06}. Whether the two notions differ is the question of footnote~56 in
its plainest form. This paper answers it and decides Corollaire~7.7.2.

\subsection{The axioms}\label{subsec:axioms}

$\Cat$ is the category of small categories and $\Fl(\Cat)$ its class of arrows. For a
functor $u\colon A\to B$ and an object $b$ of $B$, $A/b$ is the category of pairs $(a,\beta)$
with $a$ in $A$ and $\beta\colon u(a)\to b$. For a functor $p\colon X\to S$ and $s$ in $S$
we write $X/s$ for the corresponding comma category, and for a functor $u\colon X\to Y$ over
$S$ we write $u/s\colon X/s\to Y/s$. $BA$ denotes the classifying space $|NA|$, and $\Winf$
the class of functors $u$ such that $Bu$ is a weak homotopy equivalence. Let
$W\subseteq\Fl(\Cat)$. A small category $A$ is \emph{$W$-aspherical} if $A\to e$ lies in
$W$, where $e$ is the final category. The axioms are those of \cite[XVI]{Der}.
\begin{itemize}[leftmargin=5.8em,labelwidth=5.3em,labelsep=0.5em,align=left]
\item[Loc (1)] \emph{(Moderate saturation.)} (a) Isomorphisms lie in $W$; (b) if two of
$f$, $g$, $gf$ lie in $W$, so does the third; (c) if $gf$ and $fg$ lie in $W$, so do $f$ and
$g$.
\item[Loc (1 bis)] Loc~(1), and $W$ is closed under retracts in the category of arrows.
\item[Loc (1 ter)] \emph{(Strong saturation.)} A functor lies in $W$ as soon as its image in
the localization $\Cat[W^{-1}]$ is invertible.
\item[Loc (2)] Every small category with a final object is $W$-aspherical.
\item[Loc (3)] \emph{(Quillen's Theorem~A, absolute form.)} If $u\colon A\to B$ is such that $A/b$ is
$W$-aspherical for every object $b$ of $B$, then $u\in W$.
\item[Loc (3 bis)] \emph{(Quillen's Theorem~A, relative form.)} If $u\colon X\to Y$ is a functor over a small
category $S$ and $u/s\in W$ for every object $s$ of $S$, then $u\in W$.
\end{itemize}
Loc~(1~ter) implies Loc~(1~bis), which implies Loc~(1). A class is \emph{weakly saturated}
\cite[1.4.12]{Cis06} if it contains the identities, satisfies two out of three, and contains
$i$ (hence $r$) whenever $ri=1$ and $ir\in W$; by Loc~(1)(c) applied to $f=i$, $g=r$, every
class satisfying Loc~(1) is weakly saturated.

A \emph{Mayer--Vietoris square}, or \emph{MV square}, is a small category $X$ with two full
subcategories $X_1$, $X_2$ that are sieves (an arrow whose target lies in $X_k$ has its
source in $X_k$) and whose objects together exhaust $X$; we put $X_0=X_1\cap X_2$.
\begin{itemize}[leftmargin=5.8em,labelwidth=5.3em,labelsep=0.5em,align=left]
\item[Loc ($6'$a)] \emph{(Gluing.)} For every MV square, if the inclusion $X_0\to X_1$ lies in
$W$, so does the inclusion $X_2\to X$.
\item[Loc (8)] \emph{(Filtered colimits.)} For every small filtered category $I$ and every
functor $I\to\Fl(\Cat)$, $i\mapsto(u_i\colon A_i\to B_i)$, with all $u_i$ in $W$, the
functor $\colim_i A_i\to\colim_i B_i$ lies in $W$ (for any choice of the colimits). \emph{Loc (8 bis)} is the case
$I=\N$.
\item[Loc (9)] \emph{(Accessibility.)} $W$ is an accessible full subcategory of the category of
arrows of $\Cat$.
\end{itemize}
Grothendieck's MV squares are formed with \emph{ouverts} \cite[XVI, pp.~12--13]{Der}. Under Loc~(1),
(2), (3) one has $f\in W$ if and only if $f^{\circ}\in W$, where $f^{\circ}$ is the functor
between opposite categories \cite[XVI, Rappel~1]{Der}, and the same holds for every weak
fundamental localizer \cite[Prop.~1.1.21]{Mal05}; since passing to opposite categories
exchanges sieves and cosieves, Loc~($6'$a) has the same meaning in either convention for every
weak fundamental localizer. We shall verify Loc~(9) in the following form: for $\kappa=\aleph_1$,
$W$ is closed under $\kappa$-filtered colimits in $\Fl(\Cat)$, and every member of $W$ is the
$\kappa$-filtered union of its sub-functors that lie in $W$ and are functors between
countable categories. Since a functor between countable categories is an $\aleph_1$-presentable
object of $\Fl(\Cat)$, hence of $W$, and there is only a set of them up to isomorphism, this
makes $W$ an $\aleph_1$-accessible full subcategory, closed under $\aleph_1$-filtered colimits.
This is what Loc~(9) asks. By \cite[XVI, 5.2]{Der} a part of $\Fl(\Cat)$ is accessible in the
sense of Chapter~XVIII of the manuscript: it is closed under isomorphism, the full subcategory
it defines is accessible, that is, equivalent to a category of ind-objects of a small category,
and its inclusion is an accessible functor \cite[D\'ef.~4.13.1 and~4.8.7]{DerXVIII}. For a
class satisfying Loc~(8), the existence of the countable sub-functors is the criterion of
\cite[XVI, Cor.~5.3.2]{Der} for $\pi=\aleph_0$.

\begin{definition}\label{def:localizers}
A \emph{weak fundamental localizer} is a weakly saturated class $W\subseteq\Fl(\Cat)$
satisfying Loc~(2) and Loc~(3) \cite[1.1.2]{Mal05}, \cite[6.1.12]{Cis06}. A
\emph{fundamental localizer} is a weakly saturated class satisfying Loc~(2) and Loc~(3~bis)
\cite[D\'ef.~3.3.2]{Cis06}.
\end{definition}

Taking $Y=S$ in Loc~(3~bis) gives Loc~(3) (the comma categories $S/s$ have final objects), so
every fundamental localizer is a weak fundamental localizer \cite[6.1.12]{Cis06}. A localiseur
fondamental in the sense of \cite[XVI, D\'ef.~1.1]{Der} is a weak fundamental localizer that
satisfies Loc~(1).

\subsection{Results}\label{subsec:results}

For a functor $u\colon A\to B$ consider the following conditions.
\begin{itemize}[leftmargin=3.6em,labelwidth=3.1em,labelsep=0.5em,align=left]
\item[($\pi$)] $\pi_0(u)$ is bijective, and $u\in\Winf$ if $A$ is connected.
\item[(circ)] $u\in\Winf$, or $BA$ and $BB$ are homotopy equivalent to the circle and
$H_1(Bu;\Q)$ is an isomorphism.
\item[(tor)] $\Pione(Bu)$ is an equivalence of fundamental groupoids, $H_n(Bu;\Z)$ is an
isomorphism for every $n\ge3$, $H_2(Bu;\Q)$ is an isomorphism, and $Bu$ induces an
isomorphism of the torsion subgroups $\Tor H_2(BA;\Z)\to\Tor H_2(BB;\Z)$.
\item[(fin)] $\Pione(Bu)$ is an equivalence, $H_*(Bu;\Z)$ is an isomorphism, and $Bu$
restricts to a weak homotopy equivalence on all but finitely many connected components of
$BA$.
\end{itemize}
Let $\Wpi$, $\Wcirc$, $\Wtor$, $\Wfin$ be the classes of functors between small categories
satisfying ($\pi$), (circ), (tor), (fin) respectively.

\begin{theoremA}
Each of the classes $\Wpi$, $\Wcirc$, $\Wtor$, $\Wfin$ satisfies Loc~\textup{(1)},
\textup{(2)} and \textup{(3)}; it is therefore a weak fundamental localizer, and a localiseur
fondamental in the sense of \cite[XVI, D\'ef.~1.1]{Der}. None of them satisfies
Loc~\textup{(3~bis)}, so none is a fundamental localizer. For each of them the aspherical
categories are exactly the weakly contractible categories. Moreover:
\begin{enumerate}
\item[(a)] $\Wpi$ satisfies neither Loc~\textup{(1~bis)}, nor Loc~\textup{($6'$a)}, nor
Loc~\textup{(8~bis)}, nor Loc~\textup{(9)}. Loc~\textup{(3~bis)} fails for the projection $E\times BC_2\to E$ over
$BC_2$, where $E$ is the groupoid with two objects and exactly one arrow between any two
objects, mapped to $BC_2$ as the translation groupoid of the group $C_2$ of order two.
\item[(b)] $\Wcirc$ is strongly saturated and satisfies neither Loc~\textup{($6'$a)} nor
Loc~\textup{(8~bis)}.
Loc~\textup{(3~bis)} fails over the ordered set $\{o<+,\ o<-\}$, for an order-preserving map
from an ordered set with $10$ elements to one with $6$.
\item[(c)] $\Wtor$ is strongly saturated and satisfies Loc~\textup{(8)} and
Loc~\textup{(9)}; it does not satisfy Loc~\textup{($6'$a)}. Loc~\textup{(3~bis)} fails over the
ordered set $\{0<1,\ 0<2\}$, for an order-preserving map from an ordered set with $77$
elements to one with $53$.
\item[(d)] $\Wfin$ is strongly saturated and satisfies Loc~\textup{($6'$a)} and
Loc~\textup{(9)}; it does not satisfy Loc~\textup{(8~bis)}. Loc~\textup{(3~bis)} fails over the
discrete category $\N$, for a coproduct of copies of one functor between finite categories
with $2$ and $58$ objects.
\end{enumerate}
\end{theoremA}

\begin{theoremB}
Let $W$ be a weak fundamental localizer.
\begin{enumerate}
\item[(i)] If $W$ satisfies Loc~\textup{($6'$a)}, then $W$ is strongly saturated.
\item[(ii)] $W$ is a fundamental localizer if and only if it satisfies Loc~\textup{($6'$a)}
and Loc~\textup{(8)}.
\end{enumerate}
\end{theoremB}

\begin{corollaryC}
Let $W$ be a class of functors satisfying Loc~\textup{(1)}, \textup{(2)}, \textup{(3)}.
\begin{enumerate}
\item[(i)] Grothendieck's Corollaire~7.7.2, $\mathrm{Loc}\,(6'\mathrm{a},8,9)\Rightarrow
\mathrm{Loc}\,(1\ \mathrm{ter},3\ \mathrm{bis})$, holds. More precisely Loc~\textup{($6'$a)}
alone gives Loc~\textup{(1~ter)}, and Loc~\textup{($6'$a)} with Loc~\textup{(8)} gives
Loc~\textup{(3~bis)}; Loc~\textup{(9)} is not needed.
\item[(ii)] Neither Loc~\textup{($6'$a)} nor Loc~\textup{(8)} can be omitted from the
hypotheses giving Loc~\textup{(3~bis)}, even if Loc~\textup{(1~ter)} and Loc~\textup{(9)} are
kept: $\Wtor$ and $\Wfin$ are counterexamples. Nor can Loc~\textup{(8)} be weakened to
Loc~\textup{(8~bis)}: the class $\Wctbl$ of Proposition~\ref{prop:ctbl} is a counterexample.
\item[(iii)] Loc~\textup{(1)}, \textup{(2)}, \textup{(3)} alone imply neither
Loc~\textup{(1~bis)} nor Loc~\textup{(3~bis)}.
\end{enumerate}
\end{corollaryC}

Table~\ref{tab:classes} collects the properties of the four classes and of the class
$\Wctbl$ of Proposition~\ref{prop:ctbl}. Every class that is not a fundamental localizer must,
by Theorem~B, fail Loc~($6'$a) or Loc~(8), and each of the five does.

\begin{table}[ht]
\centering
\small
\begin{tabular}{@{}lllllll@{}}
\toprule
Class & Saturation & ($6'$a) & (8) & (9) & Base of the witness & Size of the witness\\
\midrule
$\Wpi$ & Loc (1), not (1 bis) & no & no & no & $BC_2$ & groupoids, $2\to2$ objects\\
$\Wcirc$ & Loc (1 ter) & no & no & -- & $\{o<+,\ o<-\}$ & posets, $10\to6$\\
$\Wtor$ & Loc (1 ter) & no & yes & yes & $\{0<1,\ 0<2\}$ & posets, $77\to53$\\
$\Wfin$ & Loc (1 ter) & yes & no & yes & discrete $\N$ & $\coprod_{\N}$ of $2\to58$ objects\\
$\Wctbl$ & Loc (1 ter) & yes & no$^{*}$ & -- & discrete, uncountable & $\coprod$ of $2\to58$ objects\\
\bottomrule
\end{tabular}
\caption{The four weak fundamental localizers of Theorem~A and the class $\Wctbl$ of
Proposition~\ref{prop:ctbl}. A dash means that the property is not decided here. Each class
fails Loc~(3~bis). ${}^{*}$\,$\Wctbl$ satisfies Loc~(8~bis); $\Wpi$ and $\Wcirc$ fail it.}
\label{tab:classes}
\end{table}

\subsection{What the results mean}\label{subsec:meaning}

Theorem~A shows, for the literal axioms of Chapter~XVI and for every strengthening of the
saturation axiom that Grothendieck considers, that the absolute form of Quillen's Theorem~A, Loc~(3), does
not imply the relative form, Loc~(3~bis). Theorem~B says exactly what
has to be added: the gluing axiom and stability under filtered colimits. Both are properties
of every fundamental localizer (Proposition~\ref{prop:necessity}), so for a weak fundamental
localizer, being a fundamental localizer \emph{is} the conjunction of Loc~($6'$a) and
Loc~(8). Corollary~C settles Corollaire~7.7.2: it is true as stated and remains true without
the accessibility axiom, while the doubt of footnote~56 is borne out as soon as Loc~($6'$a)
or Loc~(8) is dropped.

\subsection{Why it matters}\label{subsec:why}

The axioms of fundamental localizers are the foundation of Grothendieck's homotopy theory of
small categories, as developed by Maltsiniotis \cite{Mal05} and Cisinski \cite{Cis06}: test
categories, the comparison of homotopy theories of presheaves, and the derivators
$\mathrm{HOT}_W$ of Chapter~XVI all rest on them. Since much of the theory needs only the
absolute form, the weak notion was isolated and shown to agree with the strong one at the
bottom, for $\Winf$ \cite[Th.~6.1.18]{Cis06}, and for the weak fundamental localizers
generated by the aspherical categories of a fundamental localizer \cite[Cor.~6.1.17]{Cis06}.
Theorem~A shows that the two notions differ in general, by finite examples (for $\Wfin$, a
coproduct of copies of one finite example), and that the
difference cannot be seen on aspherical categories: our classes have the same aspherical
categories as $\Winf$ and are strictly larger. A class of weak equivalences defined by
homological or homotopical invariants, which one wants to use as a localizer, therefore has to
be checked against the relative axiom, or against Loc~($6'$a) and Loc~(8); Theorem~B makes
that check practical, and once it is passed Cisinski's theory applies in full. For the reading
of \emph{Les D\'erivateurs}, the corollary removes a doubt that Grothendieck left in his text:
the \emph{localiseurs sp\'eciaux} of XVI~7.8, those satisfying Loc~($6'$a,~8,~9), satisfy
Loc~(3~bis): his footnote~57 hesitates to add Loc~(3~bis) to their definition, and footnote~58
notes that recovering $W$ from $\mathrm{HOT}_W$ needs it; the axiom that both footnotes ask for
holds automatically.

\subsection{Credit}\label{subsec:credit}

The axioms Loc~(1)--(9), the question, Corollaire~7.7.2 and its footnotes are Grothendieck's
\cite[XVI]{Der}; his footnote~53 \cite[XVI, p.~53]{Der} announces the implication
``Loc (6 a, 8, 9) implique d\'ej\`a Loc (3 bis)'' that footnote~56 later doubts. The notion of
weak fundamental localizer is Maltsiniotis's \cite[1.1.2]{Mal05}. The formal definition of
fundamental localizers in the relative form, the minimality theorem \cite[Th.~6.1.18]{Cis06},
the theory of localizers of simplicial sets and its correspondence with fundamental
localizers \cite[Th.~4.2.15]{Cis06}, and the stability of fundamental localizers under
filtered colimits \cite[Cor.~4.2.22]{Cis06} are Cisinski's; the proof of Theorem~B rests on
them. Cisinski also proves that every fundamental localizer is strongly saturated
\cite[Prop.~4.2.4]{Cis06}; Theorem~B(i) extends this to the weak fundamental localizers that
satisfy Loc~($6'$a). The saturation criterion used for Theorem~B(i) is Rădulescu-Banu's
\cite[Th.~7.2.7]{RB09}, and Quillen's Theorem~A \cite[\S1]{Qui73} is used throughout. Supplied
here are the four classes and their witnesses (Theorem~A), the sufficiency in Theorem~B(ii)
and Theorem~B(i) as deductions from the results just named, and the resulting decision of
Corollaire~7.7.2. Proposition~\ref{prop:ctbl}, the failure of Loc~(8~bis) for $\Wcirc$ and the
failure of Loc~(9) for $\Wpi$ are due to the checker of this paper. Section~\ref{sec:method}
says which AI model found which part.

\subsection{Context}\label{subsec:context}

Cisinski introduces weak fundamental localizers in \cite[6.1.12]{Cis06} and proves three
comparison results: a weak fundamental localizer contains every proper fundamental localizer
whose aspherical categories it makes aspherical \cite[Prop.~6.1.16]{Cis06}; the smallest weak
fundamental localizer making the aspherical categories of a fundamental localizer $W$
aspherical is a fundamental localizer \cite[Cor.~6.1.17]{Cis06}; and $\Winf$ is the smallest
weak fundamental localizer \cite[Th.~6.1.18]{Cis06}. All three are statements about the weak
fundamental localizers generated by asphericity. Theorem~A concerns the others: each of our
classes has the same aspherical categories as $\Winf$, the smallest weak fundamental localizer
\cite[Th.~6.1.18]{Cis06}, and is strictly larger than $\Winf$. In the other direction
Cisinski proves that every fundamental localizer is stable under filtered colimits
\cite[Cor.~4.2.22]{Cis06}; Proposition~\ref{prop:necessity} adds the gluing axiom by the same
method, and Theorem~B shows that the two properties together characterize fundamental
localizers among weak ones. Grothendieck's own route to Loc~(3~bis) in XVI~7.7 passes through
the derivator $\mathrm{HOT}_W$; ours passes through localizers of simplicial sets and needs
neither a model structure on $\Cat$ nor a derivator, which is how it avoids the circle of
footnote~55. We know of no earlier example of a weak fundamental localizer that is not a
fundamental localizer; Maltsiniotis's and Cisinski's treatments \cite[\S1.1]{Mal05},
\cite[\S6.1]{Cis06} give none.

\subsection{Scope}\label{subsec:scope}

What follows: Corollaire~7.7.2 and its sharpening (Corollary~C); the localiseurs sp\'eciaux of
XVI~7.8 satisfy Loc~(3~bis); a weak fundamental localizer that satisfies Loc~(8~bis) is closed
under retracts (Lemma~\ref{lem:retracts}). What is not proved: whether Loc~(8) can be replaced
in Theorem~B by stability under coproducts (Loc~(8~bis) does not suffice,
Proposition~\ref{prop:ctbl}); whether Loc~(8) and Loc~(9)
without Loc~($6'$a) imply strong saturation; the entries marked by a dash in
Table~\ref{tab:classes}. These are stated in Section~\ref{sec:open}. We do not discuss the
other axioms of Chapter~XVI (Loc~(4), (5), (6~b), (7)), nor derivators.

\section{Proof overview}\label{sec:overview}

The absolute form of Quillen's Theorem~A, Loc~(3), constrains a class $W\supseteq\Winf$ only through its
aspherical categories. If these are exactly the weakly contractible categories, Quillen's
Theorem~A does the rest: a functor with weakly contractible comma categories lies in $\Winf$,
hence in $W$ (Lemma~\ref{lem:absolute}). The relative form is of another nature: it is a
gluing property. For a functor $D\colon J\to\Cat$ the comma categories of the Grothendieck
construction $\int_J D\to J$ recover the values $D(j)$ up to adjunction
(Lemma~\ref{lem:grothendieck}), and the classifying space of $\int_J D$ is the homotopy colimit
of the spaces $BD(j)$ \cite{Tho79}. So Loc~(3~bis) asks that $W$ be closed under homotopy
colimits of diagrams of categories, and a separating class is a class that contains $\Winf$,
has only weakly contractible aspherical categories, and is not closed under one kind of
homotopy colimit. Theorem~A gives four:
\begin{itemize}
\item[(a)] $\Wpi$ forgets everything about a category with several components except their
number; the homotopy orbit space of a free action of $C_2$ on two points is a single point, and
over $BC_2$ two components become one.
\item[(b)] $\Wcirc$ accepts maps of circles of non-zero degree but not their suspensions; over
the three-element ordered set $\{o<+,\ o<-\}$ the homotopy colimit is a suspension.
\item[(c)] $\Wtor$ accepts the map of degree $2$ of the $2$-sphere, which is a rational
isomorphism and creates no torsion; its cofibre, a homotopy pushout over $\{0<1,\ 0<2\}$, is
the Moore space with $H_2=\Z/2$.
\item[(d)] $\Wfin$ accepts any map that fails to be a weak equivalence on finitely many
components; over the discrete category $\N$ the homotopy colimit is an infinite coproduct.
\end{itemize}
The finite witness of (d) is a functor whose classifying space is the inclusion of a circle
into the mapping torus of a self-map of degree $2$ of the $2$-sphere: it induces isomorphisms on the fundamental group and on integral
homology, but the universal cover of the mapping torus has $H_2=\Z[1/2]$.

For the positive direction, Theorem~B(ii), let
$V=\{f\colon X\to Y\ \text{in}\ \sSet: \Delta/f\in W\}$, where $\Delta/X$ is the category of
simplices. Cisinski's correspondence \cite[Th.~4.2.15]{Cis06} turns $V$ into a fundamental
localizer on $\Cat$ as soon as $V$ is a localizer of simplicial sets in his sense
\cite[D\'ef.~1.4.1]{Cis06}: two out of three, trivial fibrations in $V$, and monomorphisms in
$V$ stable under pushouts and transfinite composition. The fundamental localizer obtained is
$W$ itself, because the functor $\Delta/NA\to A$ is a natural weak equivalence. Two out of
three and the trivial fibrations come from $W\supseteq\Winf$ \cite[Th.~6.1.18]{Cis06}. The
pushouts are where Loc~($6'$a) is used: a pushout of simplicial sets along a monomorphism
becomes, after replacing the other map by its mapping cylinder, a union of two simplicial
subsets, and the categories of simplices of two simplicial subsets form an MV square
(Lemma~\ref{lem:pushout}). The transfinite compositions are where Loc~(8) is used, at limit
ordinals. Loc~(9) is used nowhere. Theorem~B(i) uses only the pushout step, together with a
saturation theorem for cofibration categories \cite[Th.~7.2.7]{RB09}. The point of this route
is that the relative axiom is never assumed on the way; Grothendieck's argument in XVI~7.7
identified $W$ with the localizer of the derivator $\mathrm{HOT}_W$, which he had proved under
Loc~(3~bis).

\section{Preliminaries}\label{sec:prelim}

This section gives the tools shared by the four examples and the proof of Theorem~B.

\begin{lemma}\label{lem:absolute}
Let $W$ be a weakly saturated class with $\Winf\subseteq W$ whose aspherical categories are
weakly contractible. Then $W$ satisfies Loc~\textup{(2)} and Loc~\textup{(3)}, so it is a weak
fundamental localizer, and its aspherical categories are exactly the weakly contractible
categories.
\end{lemma}

\begin{proof}
A category with a final object has a contractible classifying space \cite[\S1, Corollary~2 of Proposition~2]{Qui73},
so $A\to e$ lies in $\Winf\subseteq W$; this is Loc~(2), and the same argument shows that
every weakly contractible category is $W$-aspherical. If $u\colon A\to B$ has $W$-aspherical
comma categories $A/b$, these are weakly contractible by hypothesis, and Quillen's Theorem~A
\cite[\S1, Th.~A]{Qui73}, in its dual formulation for the comma categories $A/b$, which Quillen
states right after the theorem, shows that
$u\in\Winf\subseteq W$. This is Loc~(3).
\end{proof}

\begin{lemma}\label{lem:inverse-image}
Let $F\colon\Cat\to\mathcal D$ be a functor to any category, and let $W$ be the class of
functors $u$ such that $F(u)$ is invertible. Then $W$ is strongly saturated; in particular it
satisfies Loc~\textup{(1~bis)} and is weakly saturated.
\end{lemma}

\begin{proof}
$F$ inverts $W$, so it factors as $F=\bar F\gamma$ through the localization
$\gamma\colon\Cat\to\Cat[W^{-1}]$ (formed if necessary in a larger universe). If $\gamma(u)$
is invertible, so is $F(u)=\bar F(\gamma(u))$, and $u\in W$. A strongly saturated class
satisfies Loc~(1~bis), because $\gamma$ sends a retract of an arrow with invertible image to
an arrow with invertible image.
\end{proof}

For a functor $D\colon J\to\Cat$ the \emph{Grothendieck construction} $\int_JD$ has as
objects the pairs $(j,x)$ with $x$ in $D(j)$, and as arrows $(j,x)\to(j',x')$ the pairs
$(\alpha,\varphi)$ with $\alpha\colon j\to j'$ and $\varphi\colon D(\alpha)x\to x'$; it comes
with the projection to $J$. A natural transformation $\phi\colon D\to D'$ induces a functor
$\int\phi\colon\int_JD\to\int_JD'$ over $J$.

\begin{lemma}\label{lem:grothendieck}
For every object $j$ of $J$ the functor
$\mathrm{ev}_j\colon(\int_JD)/j\to D(j)$, $\bigl((j',x),\alpha\colon j'\to j\bigr)\mapsto
D(\alpha)x$, is left adjoint to $x\mapsto((j,x),\id_j)$, and $\mathrm{ev}_j\circ(\int\phi)/j=
\phi_j\circ\mathrm{ev}_j$. Consequently, if $W\supseteq\Winf$ satisfies two out of three, then
$(\int\phi)/j\in W$ if and only if $\phi_j\in W$.
\end{lemma}

\begin{proof}
An arrow from $((j',x),\alpha)$ to $((j,y),\id_j)$ in $(\int_JD)/j$ is an arrow
$(\gamma,\psi)$ of $\int_JD$ with $\id_j\circ\gamma=\alpha$; it is therefore determined by an
arrow $\psi\colon D(\alpha)x\to y$ of $D(j)$, naturally in both variables. The second
formula is immediate. A functor with an adjoint lies in $\Winf$ \cite[\S1, Corollary~1 of Proposition~2]{Qui73}, and
the last assertion follows by two out of three.
\end{proof}

Two special cases recur. For a functor $f\colon A\to B$, the \emph{collage} $M_f$ is the
Grothendieck construction over the ordered set $\{0<1\}$ of the diagram $f$: its objects are
those of $A$ and of $B$, with the arrows of $A$ and $B$, arrows $a\to b$ given by the arrows
$f(a)\to b$ of $B$, and no arrows from $B$ to $A$. $A$ is a sieve in $M_f$. The functor
$r\colon M_f\to B$ that is $f$ on $A$ and the identity on $B$ is left adjoint to the inclusion
of $B$, and $r$ restricted to $A$ is $f$; hence for $W\supseteq\Winf$ with two out of three,
the inclusion $A\to M_f$ lies in $W$ if and only if $f$ does. For a span
$B\xleftarrow{f}A\xrightarrow{g}C$, the \emph{double collage} is the Grothendieck construction
over $J=\{0<1,\ 0<2\}$; it is the union of $M_f$ and $M_g$ along $A$, with no arrows between
$B$ and $C$.

\begin{lemma}\label{lem:mv}
For an MV square $X=X_1\cup X_2$, the nerve $NX$ is the union of the simplicial subsets
$NX_1$ and $NX_2$, with intersection $NX_0$. Consequently $BX$ is the homotopy pushout of
$BX_1\leftarrow BX_0\to BX_2$, and $\pi_0(BX)$ is the pushout of the sets $\pi_0(BX_k)$. In a
double collage $P$ of $B\xleftarrow{f}A\xrightarrow{g}C$, the collages $M_f$ and $M_g$ form an
MV square with intersection $A$, and $BP$ is homotopy equivalent to the homotopy pushout of
$BB\xleftarrow{Bf}BA\xrightarrow{Bg}BC$.
\end{lemma}

\begin{proof}
A simplex of $NX$ is a chain $x_0\to\dots\to x_n$; its last object lies in $X_1$ or $X_2$, and
since these are sieves and full, the whole chain lies there. A chain lying in both lies in
$X_0$. The realization of a union of two simplicial subsets is a union of two subcomplexes, hence
a homotopy pushout. For the double collage, the full subcategories $M_f$ (objects of $A$ and
$B$) and $M_g$ (objects of $A$ and $C$) are sieves, they exhaust $P$, and they meet in $A$. The
retractions $r\colon M_f\to B$ and $M_g\to C$ are weak equivalences restricting to $f$ and
$g$ on $A$, so the homotopy pushout of $BM_f\leftarrow BA\to BM_g$ is that of
$BB\leftarrow BA\to BC$.
\end{proof}

\begin{lemma}\label{lem:retracts}
A class $W$ satisfying Loc~\textup{(8~bis)} is closed under retracts in the category of
arrows of $\Cat$.
\end{lemma}

\begin{proof}
Let $f$ be a retract of $g\in W$ in the category of arrows, with $s\colon f\to g$ and
$r\colon g\to f$ such that $rs=\id_f$, and let $e=sr$. The sequence
$g\xrightarrow{e}g\xrightarrow{e}g\to\cdots$ has colimit $f$, with the maps $r$ as colimit
cone: $re=r$; and if maps $c_n\colon g\to Z$ satisfy $c_{n+1}e=c_n$, then
$c_ne=c_{n+1}ee=c_{n+1}e=c_n$, so all $c_n$ are equal to $c_0$, and $c_0s$ is the unique map
$\varphi$ with $\varphi r=c_0$ (indeed $c_0sr=c_0e=c_0$, and $\varphi=\varphi rs$). Loc~(8~bis)
gives $f\in W$.
\end{proof}

\section{Four separating classes}\label{sec:examples}

This section proves Theorem~A, one class at a time; for the strong saturation of $\Wfin$ it
uses Theorem~B(i), proved in Section~\ref{sec:positive} independently of this section.

\subsection{Components: the class \texorpdfstring{$\Wpi$}{W pi}}\label{subsec:wpi}

\emph{Loc (1).} Isomorphisms satisfy ($\pi$). If two of $f$, $g$, $gf$ satisfy ($\pi$), then
all three are bijective on components, so the three categories involved have the same number
of components. If that number is not $1$ nothing more is required; if it is $1$, two of the
three functors lie in $\Winf$, and so does the third. If $gf$ and $fg$ satisfy ($\pi$), then
$\pi_0(f)$ and $\pi_0(g)$ are bijective; if the categories are connected, $Bf$ has a left
inverse $(Bgf)^{-1}Bg$ and a right inverse $Bg(Bfg)^{-1}$ in the homotopy category of spaces,
so $Bf$ is a homotopy equivalence of CW complexes; symmetrically so is $Bg$, and $f,g\in\Wpi$.

\emph{Loc (2), (3).} A functor $A\to e$ satisfies ($\pi$) exactly when $A$ is connected and
$BA$ is weakly contractible. Since $\Winf\subseteq\Wpi$, Lemma~\ref{lem:absolute} applies.

\emph{Loc (3 bis) fails.} Write $G=C_2=\{0,1\}$ additively, and let $C=BG$ be the category with
one object $c$ and endomorphisms $G$. Let $E$ be the groupoid with objects $x\in G$ and one
arrow $(x,g)\colon x\to x+g$ for each $g\in G$, composed by
$(x+g,h)\circ(x,g)=(x,g+h)$, and let $p\colon E\to C$ send $(x,g)$ to $g$. Let
$A=E\times BG$, let $u\colon A\to E$ be the projection, and regard $u$ as a functor over $C$
through $pu$ and $p$. The groupoid $E$ is connected with trivial automorphism groups, so $BE$
is contractible; $A$ is connected with automorphism group $G$, so $BA\simeq BG$ has
fundamental group $\Z/2$. Hence $u\notin\Winf$, and since $A$ is connected, $u\notin\Wpi$.

The comma category $E/c$ has as objects the pairs $(x,t)\in G\times G$, with $t\colon p(x)\to
c$, and an arrow $(x,g)$ from $(x,t)$ to $(x+g,t')$ exactly when $t'+g=t$. The sum $x+t$ is
invariant under arrows, and each of its two values is taken by two objects joined by exactly
one arrow in each direction. So $E/c$ has exactly two components, each equivalent to $e$. Since
$pu$ ignores the factor $BG$, $A/c=(E/c)\times BG$, and $u/c$ is the projection; both
categories have two components and $u/c$ is bijective on them. As $A/c$ is not connected,
$u/c\in\Wpi$. This is the only object of $C$, so the premise of Loc~(3~bis) holds while its
conclusion fails.

\emph{Loc (1 bis), (8 bis) and ($6'$a) fail.} The functor $BG\to e$ is a retract, in the
category of arrows, of $BG\amalg e\to e\amalg e$ (collapse on the first summand, identity on
the second), using the inclusions of the first summands and the retractions sending the
added object to the unique object; the latter functor lies in $\Wpi$ because its source has two
components, the former does not. By Lemma~\ref{lem:retracts}, $\Wpi$ does not satisfy
Loc~(8~bis). For Loc~($6'$a), let $X$ be the ordered set with elements $0,1,2,3,q,m$, where
$0<1$, $0<3$, $2<1$, $2<3$ (a square whose nerve is a circle), $0<m$ and $q<m$. The subsets
$X_1=\{0,1,2,3,q\}$ and $X_2=\{0,q,m\}$ are sieves exhausting $X$, with $X_0=\{0,q\}$. The
inclusion $X_0\to X_1$ is bijective on components, and $X_0$ has two of them, so it lies in
$\Wpi$; but $X_2$ is connected and contractible while $BX$ is a circle with an arc attached,
so $X_2\to X$ is not in $\Wpi$.

\emph{Loc (9) fails.} By \cite[XVI, 5.2]{Der} and \cite[D\'ef.~4.13.1 and~4.8.7]{DerXVIII}, an
accessible part of $\Fl(\Cat)$ defines a full subcategory equivalent to a category of
ind-objects $\mathrm{Ind}_{\pi_0}(C)$ of a small category $C$, for some infinite cardinal
$\pi_0$; such a category has colimits of all diagrams indexed by ordered sets in which every
subset of cardinality at most $\pi_0$ has an upper bound. We show that the full subcategory
$\Wpi$ of the category of arrows lacks such colimits, whatever $\pi_0$. Let $\kappa>\pi_0$ be a
regular cardinal, $\Omega$ a set of cardinality $\kappa$, and $I$ the set of subsets of $\Omega$
of cardinality less than $\kappa$, ordered by inclusion; a subset of $I$ of cardinality at most
$\pi_0$ is bounded by its union. For $i\in I$ let $A_i=BG\amalg(\Omega\smallsetminus i)$ and
$B_i=\{\ast\}\amalg(\Omega\smallsetminus i)$, with $\Omega\smallsetminus i$ discrete, and let
$u_i\colon A_i\to B_i$ collapse $BG$ to $\ast$ and be the identity elsewhere; for $i\subseteq j$
the transition functors send the points of $j\smallsetminus i$ to the object of $BG$ and to
$\ast$. Each $u_i$ is bijective on components and its source is not connected, so
$u_i\in\Wpi$; the colimit of the diagram in $\Fl(\Cat)$ is $c\colon BG\to e$. A colimit
$w\colon P\to R$ of the diagram in the full subcategory $\Wpi$ would be a reflection of $c$:
an arrow $\lambda=(\alpha,\beta)$ from $c$ to $w$ through which every arrow from $c$ to a member
of $\Wpi$ factors uniquely. Then $P$ is connected: otherwise let $P_0$ be the component
containing $\alpha(BG)$, $P_1\ne\varnothing$ the rest, and $R_0$, $R_1$ the corresponding parts
of $R$; the functor $w\amalg\id_e$ lies in $\Wpi$, and the inclusion of $w$ into it and the
arrow that is the identity on $P_0$, $R_0$ and sends $P_1$, $R_1$ to the added object are two
different arrows that agree on $\lambda$, since $\beta$ lands in $R_0$. So $w\in\Winf$. The
arrow from $c$ to the member $BG\amalg e\to e\amalg e$ of $\Wpi$ given by the inclusions of
the first summands factors through $\lambda$, which gives a functor $\nu\colon P\to BG\amalg e$
with $\nu\alpha$ the inclusion of $BG$; as $P$ is connected, $\nu$ is a retraction $P\to BG$ of
$\alpha$. But $w\alpha$ factors through $e$ and $Bw$ is a homotopy equivalence, so $B\alpha$ is
null-homotopic, and then so is the identity of $B(BG)$, the infinite real projective space,
which is absurd. So $\Wpi$ is not an accessible part of $\Fl(\Cat)$.

\subsection{Circles: the class \texorpdfstring{$\Wcirc$}{W circ}}\label{subsec:wcirc}

Let $\mathcal K$ be the homotopy category of CW complexes, and let $I(X,Y)\subseteq\mathcal
K(X,Y)$ be the set of classes that factor through the circle $S^1$. Let $\mathcal Q$ be the
category with the same objects and with
\[
\mathcal Q(X,Y)=\bigl(\mathcal K(X,Y)\amalg\{0\}\bigr)\big/\bigl(I(X,Y)\sim 0\bigr),
\]
composition being that of $\mathcal K$ on classes outside $I$ and absorbing on $0$; this is
well defined and associative because $I$ is a two-sided ideal. Let $q\colon\mathcal K\to\mathcal
Q$ be the quotient functor.

\begin{lemma}\label{lem:quotient}
$q(f)$ is invertible if and only if $f$ is invertible in $\mathcal K$, or the source and target
of $f$ are both homotopy retracts of $S^1$. A homotopy retract of $S^1$ is homotopy equivalent
to a point or to $S^1$.
\end{lemma}

\begin{proof}
$q(\id_X)=0$ exactly when $\id_X$ factors through $S^1$, that is, when $X$ is a homotopy
retract of $S^1$; such an $X$ is a zero object of $\mathcal Q$, all of whose hom-sets are
singletons. If neither $q(\id_X)$ nor $q(\id_Y)$ is $0$ and $q(f)$ has an inverse, the inverse is
not $0$, so it is $q(g)$ for some $g$, and since non-zero classes of $\mathcal Q$ are
singletons, $gf=\id_X$ and $fg=\id_Y$ in $\mathcal K$. If exactly one of the identities is $0$,
no morphism is invertible; if both are, every morphism is. A homotopy retract $X$ of $S^1$ is
connected with $\pi_n(X)=0$ for $n\ge2$, and $\pi_1(X)$ is a retract of $\Z$, which is $0$ or
$\Z$; by Whitehead's theorem \cite[Th.~4.5]{Hat02}, $X$ is a point or $S^1$ up to homotopy.
\end{proof}

\emph{Saturation and Loc (2), (3).} By Lemma~\ref{lem:quotient}, a map $f$ of CW complexes is
inverted by $(q,H_1(-;\Q))\colon\mathcal K\to\mathcal Q\times\mathrm{Vect}_{\Q}$ if and only
if it is a homotopy equivalence or a map between spaces homotopy equivalent to $S^1$ inducing
an isomorphism on $H_1(-;\Q)$: between a point and a circle the rational $H_1$ test fails. So
$\Wcirc$ is the class of functors $u$ for which this functor inverts $Bu$, and it is strongly
saturated by Lemma~\ref{lem:inverse-image}. It contains $\Winf$; since $e$ is not of circle
type, $A\to e$ lies in $\Wcirc$ only when $BA$ is weakly contractible, and
Lemma~\ref{lem:absolute} gives Loc~(2) and~(3).

\emph{Loc (3 bis) fails.} For even $m\ge4$ let $C_m$ be the ordered set on $\{0,\dots,m-1\}$
whose only relations are $2i<2i-1$ and $2i<2i+1$ (indices mod $m$). Its nerve is the polygon with
$m$ edges, and the $1$-cycle $z_m=\sum_{i\ \text{even}}\bigl([i,i+1]-[i,i-1]\bigr)$ generates
$H_1(NC_m;\Z)$. The map $v\colon C_8\to C_4$, $i\mapsto i\bmod4$, preserves the order and
$v_*z_8=2z_4$, so $Bv$ has degree~$2$ and $v\in\Wcirc$. Let $\widehat C_m$ be $C_m$ with two
incomparable elements $t_+$, $t_-$ added above every element; its nerve is the union of two
cones on $NC_m$, a $2$-sphere, and the cycle obtained by coning $z_m$ to $t_+$ minus its cone
to $t_-$ generates $H_2(N\widehat C_m;\Z)$. Let $S=\{o<+,\ o<-\}$, map $C_m$ to $o$ and $t_\pm$
to $\pm$, and let $\hat v\colon\widehat C_8\to\widehat C_4$ extend $v$ by fixing $t_\pm$; it
is a map over $S$ between ordered sets with $10$ and $6$ elements. For ordered sets the comma
category over $s$ is the inverse image of $\{s'\le s\}$. Over $o$ the comparison $\hat v/o$ is
$v\in\Wcirc$; over $\pm$ it is a map between cones $C_m\cup\{t_\pm\}$, which have greatest
elements, so it lies in $\Winf$. But $\hat v$ maps the generator of $H_2$ to twice the
generator, so it is not a weak equivalence, and its source and target are $2$-spheres, not
circles: $\hat v\notin\Wcirc$.

\emph{Loc ($6'$a) fails.} In the collage $M_v$ glue a cone point $t$ above $C_8$: the ordered
set $X=M_v\cup\{t\}$ has the sieves $X_1=M_v$ and $X_2=C_8\cup\{t\}$, with $X_0=C_8$. The
inclusion $X_0\to X_1$ lies in $\Wcirc$, since its composite with the retraction $M_v\to C_4$
is $v$. By Lemma~\ref{lem:mv}, $BX$ is the cofibre of a map of degree $2$ of the circle, the
real projective plane, with $H_1=\Z/2$; so the inclusion of the contractible $X_2$ into $X$ is
not in $\Wcirc$.

\emph{Loc (8 bis) fails.} Regard the additive group $\Z$ as a category with one object, let
$d\colon\Z\to\Z$ be the functor $m\mapsto 2m$, and let $D$ be the diagram on the ordered set
$\N$ with $D(k)=\Z$ and $D(k\le l)=d^{\,l-k}$. Let $T$ be its Grothendieck construction and
$T_n\subseteq T$ the full subcategory over $\{0,\dots,n\}$, the Grothendieck construction of the
restriction of $D$; an arrow of $T_n$ is a pair $(k\le l,\ m)$ with $m\in\Z$. Over
$\{0,\dots,n\}$ the comma category at the greatest element $n$ is all of $T_n$, so by
Lemma~\ref{lem:grothendieck} the functor $r_n\colon T_n\to\Z$,
$(k\le l,\ m)\mapsto 2^{\,n-l}m$, is a left adjoint, and $BT_n$ is a circle. The inclusion
$u_n\colon T_0\to T_n$ satisfies $r_nu_n=d^{\,n}$, of degree $2^n$, so $u_n\in\Wcirc$. With
the maps $(\id_{T_0},\ T_n\subseteq T_{n+1})$ the $u_n$ form a sequence in the category of
arrows whose colimit is the inclusion $u\colon T_0\to T$. Since $r_{n+1}$ restricted to $T_n$ is
$d\,r_n$, and homotopy groups commute with filtered colimits (Section~\ref{subsec:wtor}), the
fundamental group of $BT$ is the colimit of $\Z\xrightarrow{2}\Z\xrightarrow{2}\cdots$, that is
$\Z[1/2]$, and $\pi_1(Bu)$ is the inclusion of $\Z$. So $u$ is not a weak equivalence, and $BT$
is not homotopy equivalent to a circle: $u\notin\Wcirc$, although every $u_n$ is in $\Wcirc$.

\subsection{Torsion: the class \texorpdfstring{$\Wtor$}{W tor}}\label{subsec:wtor}

\emph{Saturation and Loc (2), (3).} Each of the four conditions in (tor) asks that a functor
inverts $Bu$: the fundamental groupoid with values in the category of groupoids and
isomorphism classes of functors (where equivalences become isomorphisms), $H_n(-;\Z)$ for
$n\ge3$, $H_2(-;\Q)$, and $\Tor H_2(-;\Z)$. So $\Wtor$ is strongly saturated
(Lemma~\ref{lem:inverse-image}), and it contains $\Winf$. If $A\to e$ lies in $\Wtor$, then
$BA$ is non-empty, connected and simply connected, $H_n(BA)=0$ for $n\ge3$, and the group
$H=H_2(BA;\Z)$ satisfies $H\otimes\Q=0$ and $\Tor H=0$; the first says that $H$ is torsion, the
second then that $H=0$. So $BA$ is simply connected and acyclic, hence weakly contractible by
the Hurewicz theorem \cite[Th.~4.32]{Hat02}, and Lemma~\ref{lem:absolute} applies.

\emph{The witness.} Let $C_m$ now be the polygon with vertices $0,\dots,m-1$ and edges
$\{i,i+1\}$ (indices mod $m$), and let $L_m=C_m*\{a,b\}$ be its suspension, a simplicial
$2$-sphere. The simplicial map $h\colon L_8\to L_4$, $i\mapsto i\bmod4$, $a\mapsto a$,
$b\mapsto b$, is injective on each simplex and has degree $2$: the fundamental cycle
$z_m=\sum_{i}\bigl([a,i,i+1]-[b,i,i+1]\bigr)$ satisfies $h_*z_8=2z_4$. Let $P$ and $Q$ be the
ordered sets of non-empty simplices of $L_8$ and $L_4$, with $50$ and $26$ elements, and
$f\colon P\to Q$ the map induced by $h$. The nerves of $P$ and $Q$ are the barycentric
subdivisions of the two spheres. The map of face posets induced by a simplicial map $g$ sends the
barycentre of each simplex to the barycentre of its image. On realizations it is homotopic to
$|g|$: a point of a simplex $\sigma$ is sent by both maps into the simplex $g(\sigma)$, and the
straight-line homotopy inside $g(\sigma)$ joins them (compare the Example after Theorem~A in
\cite[\S1]{Qui73}); $g$ need not be injective on simplices. So $Bf$ has degree $2$. Since both spheres are simply connected and torsion free and multiplication by $2$
is rationally invertible, $f\in\Wtor$.

Let $J=\{0<1,\ 0<2\}$ and consider the diagrams $D=(Q\xleftarrow{f}P\to e)$ and
$D'=(Q\xleftarrow{\id}Q\to e)$ on $J$, with the natural transformation
$\phi=(f,\id_Q,\id_e)\colon D\to D'$. The Grothendieck constructions are the double
collages: $\int_JD$ is the ordered set $P\amalg Q\amalg\{t\}$ with the orders of $P$ and $Q$,
$p<q$ when $f(p)\le q$, and $p<t$ for every $p\in P$; it has $77$ elements, and $\int_JD'$,
built in the same way from $\id_Q$, has $53$. By Lemma~\ref{lem:grothendieck} the comparisons
$(\int\phi)/j$ lie in $\Wtor$ for $j=0,1,2$, because $\phi_0=f$, $\phi_1=\id_Q$ and
$\phi_2=\id_e$ do. By Lemma~\ref{lem:mv}, $B\int_JD$ is the homotopy pushout of
$BQ\xleftarrow{Bf}BP\to\ast$, the cofibre of a map of degree $2$ of the $2$-sphere, whose
reduced homology is $\Z/2$ in degree $2$ and zero otherwise; $B\int_JD'$ is the cofibre of an
identity, which is contractible. So $\int\phi$ is not an isomorphism on $\Tor H_2$, and
$\int\phi\notin\Wtor$: Loc~(3~bis) fails.

\emph{Loc ($6'$a) fails.} In $\int_JD$, the sieves $X_1=P\cup Q$ (the collage $M_f$) and
$X_2=P\cup\{t\}$ form an MV square with $X_0=P$. The inclusion $P\to M_f$ lies in $\Wtor$, as
its composite with the retraction $M_f\to Q$ is $f$; but $X_2$ has a greatest element, so it is
contractible, while $H_2(B\int_JD)=\Z/2$. So $X_2\to\int_JD$ is not in $\Wtor$.

\emph{Loc (8).} The nerve commutes with filtered colimits, an $n$-simplex being a finite chain of
arrows. The realizations of simplicial sets then form a filtered colimit in which every compact
subset, in particular every image of a sphere or of a homotopy, comes from a finite simplicial
subset and hence from some stage; so homotopy groups and fundamental groupoids commute with
filtered colimits. Homology commutes with filtered colimits, filtered colimits of abelian groups
being exact. The torsion subgroup commutes with them as well: a class that is torsion in the
colimit is represented at some stage by an element $a$ with $na=0$ already at a later stage.
Finally a filtered colimit of equivalences of groupoids is an equivalence, since hom-sets and
objects of the colimit are colimits of those of the stages. So every condition in (tor) passes to
filtered colimits, and $\Wtor$ satisfies Loc~(8).

\emph{Loc (9).} Let $u\colon A\to B$ lie in $\Wtor$, and let $A'\subseteq A$ and $B'\subseteq B$
be countable subcategories with $u(A')\subseteq B'$. The conditions in (tor) for the
restriction $u'\colon A'\to B'$ are of the following kinds: every object of $B'$ is joined by a
path to an object $u(a)$; every path in $B'$ between objects $u(a_1)$, $u(a_2)$ is homotopic to
the image of a path in $A'$; a path in $A'$ whose image is homotopic to a constant path is itself
so; every integral cycle of $NB'$ in degree $n\ge3$ is homologous to the image of a cycle of
$NA'$, and every cycle of $NA'$ whose image bounds itself bounds; the same with rational
coefficients in degree $2$; and the same for torsion classes in degree $2$, where with a source
class we also record a chain showing that a multiple of it bounds. There are countably many such
requests, since $NA'$ and $NB'$ are countable. Because $u\in\Wtor$, each request has a witness in
$A$ or $B$, and each witness, a finite zigzag, a finite homotopy of zigzags or a finite chain,
uses finitely many arrows. Adjoin all of them to $A'$ and $B'$, close under composition and under
$u$, and repeat this countably many times. The union $u''$ is a functor between countable
categories; every request concerning $u''$ involves finitely many arrows, hence already arose at
a finite stage and was answered at the next one. So $u''\in\Wtor$. Hence the countable
sub-functors of $u$ lying in $\Wtor$ form an $\aleph_1$-filtered family with union $u$, and with
Loc~(8) this is the form of Loc~(9) stated in Section~\ref{subsec:axioms}.

\subsection{Finitely many bad components: the class \texorpdfstring{$\Wfin$}{W fin}}\label{subsec:wfin}

For $u\in\Wfin$ the map $\pi_0(Bu)$ is bijective, and we call a component $c$ of $BB$
\emph{bad} if $Bu$ is not a weak equivalence from the corresponding component of $BA$ to $c$;
the condition is that there are finitely many bad components.

\emph{Loc (1).} Fundamental groupoid equivalences and homology isomorphisms satisfy the
conditions of Loc~(1). If two of $f$, $g$, $gf$ lie in $\Wfin$, identify the components of the
three categories along the bijections; on each component two out of three holds for weak
equivalences, so the bad components of the third map are among those of the other two. If
$gf$ and $fg$ lie in $\Wfin$, then for every component $c$ of the source of $f$ such that
$c$ is not bad for $gf$ and $f(c)$ is not bad for $fg$, that is, for all but finitely many $c$, the maps
$A_c\to B_{f(c)}\to A_{gf(c)}\to B_{fgf(c)}$ have weak equivalences as composites of the first
two and of the last two, so all three are weak equivalences (two out of six for weak homotopy
equivalences, which holds since they are the maps inverted in the homotopy category). Hence
$f,g\in\Wfin$.

\emph{Loc (2), (3).} $\Winf\subseteq\Wfin$. If $A\to e$ lies in $\Wfin$, then $BA$ is
non-empty, connected, simply connected and acyclic, hence weakly contractible
\cite[Th.~4.32]{Hat02}; Lemma~\ref{lem:absolute} applies.

\emph{Loc ($6'$a).} Let $X=X_1\cup X_2$ be an MV square with $i\colon X_0\to X_1$ in $\Wfin$,
and $j\colon X_2\to X$. By Lemma~\ref{lem:mv} the square of classifying spaces is a homotopy
pushout. The van Kampen theorem for fundamental groupoids \cite[\S6.7]{Bro06} shows that
$\Pione(Bj)$ is an equivalence, a homotopy pushout of groupoids along an equivalence being an
equivalence; the Mayer--Vietoris sequences and the five lemma show that $Bj$ is a homology
isomorphism. Let $\bar S$ be the image in $\pi_0(BX)$ of the finite set of bad components of
$i$. For a component $c$ of $BX$ outside $\bar S$, restrict the whole square to the parts lying
over $c$; it is still a homotopy pushout, and its map from $BX_0$ to $BX_1$ is a disjoint union
of weak equivalences, so the map from $BX_2$ to $c$ is a weak equivalence. Hence the bad
components of $j$ lie in the finite set $\bar S$, and $j\in\Wfin$. By Theorem~B(i), $\Wfin$ is
strongly saturated.

\emph{The finite category.} Let $S_n=C_n*\{N,S\}$ be the suspension of the polygon $C_n$ of
Section~\ref{subsec:wtor}, and let $L$, $K$ be the ordered sets of non-empty simplices of
$S_6$ and $S_3$, with $38$ and $20$ elements. Let $u,v\colon S_6\to S_3$ be the simplicial maps
fixing the poles and sending the vertices $0,\dots,5$ of the polygon to $0,0,1,1,2,2$ and to
$0,1,2,0,1,2$ respectively; on the polygons they have degrees $1$ and $2$, and so do their
suspensions on $H_2$. Denote also by $u,v\colon L\to K$ the induced maps, a simplex going to
its image; as in Section~\ref{subsec:wtor}, $Bu$ and $Bv$ have degrees $1$ and $2$. Let $P$ be the category with the objects of $L$ and of $K$, the arrows of $L$ and of
$K$, and, for $a\in L$ and $b\in K$, the arrows
\[
\{(0,a,b): u(a)\subseteq b\}\ \amalg\ \{(1,a,b): v(a)\subseteq b\}
\]
from $a$ to $b$, composites with arrows of $L$ or $K$ keeping the label $0$ or $1$; there are
no arrows from $K$ to $L$. This is the Grothendieck construction of the pair $u,v$ over the
category with two objects and two parallel arrows; $P$ has $58$ objects and $518$ arrows. Let
$Q\subseteq P$ be the full subcategory on the two north poles $\{N\}\in L$ and $\{N\}\in K$: it has
two objects and two parallel arrows besides the identities, and $BQ$ is a circle. Let
$j\colon Q\to P$ be the inclusion.

A chain in $P$ passes at most once from $L$ to $K$, so $NP$ is the union of the nerves of the
two collages $M_u$ (label $0$) and $M_v$ (label $1$), which meet in $NL\amalg NK$. Hence $BP$ is
the homotopy coequalizer of $Bu,Bv\colon BL\rightrightarrows BK$, two maps of $2$-spheres of
degrees $1$ and $2$; as $Bu$ is a homotopy equivalence, $BP$ is homotopy equivalent to the
mapping torus $T$ of a self-map $\sigma$ of degree $2$ of $S^2$, and $BQ$ is the circle traced
by the fixed north pole, a section of $T\to S^1$. By van Kampen's theorem $\pi_1(T)=\Z$,
generated by this circle, so $Bj$ is an isomorphism on $\pi_1$. The Mayer--Vietoris sequence of
the mapping torus, $H_n(S^2)\xrightarrow{1-\sigma_*}H_n(S^2)\to H_n(T)\to H_{n-1}(S^2)
\xrightarrow{1-\sigma_*}H_{n-1}(S^2)$, with $1-\sigma_*=-1$ on $H_2$, gives $H_*(T)\cong
H_*(S^1)$, so $Bj$ is a homology isomorphism and $j$ satisfies the first two conditions of
(fin), with one component. But the universal cover of $T$ is the infinite mapping telescope of
$\cdots\xrightarrow{\sigma}S^2\xrightarrow{\sigma}S^2\xrightarrow{\sigma}\cdots$, which is simply
connected with $H_2=\colim(\Z\xrightarrow{2}\Z\xrightarrow{2}\cdots)=\Z[1/2]$; so
$\pi_2(BP)\cong\Z[1/2]\ne0=\pi_2(BQ)$, and $j$ is not a weak equivalence. Thus $j\in\Wfin$,
with exactly one bad component.

\emph{Loc (3 bis) and Loc (8 bis) fail.} Over the discrete category $\N$ consider
$\coprod_{n\in\N}j\colon\coprod_{n}Q\to\coprod_nP$, the $n$-th summands lying over $n$. The
comparison over $n$ is $j\in\Wfin$, but $\coprod_nj$ has infinitely many bad components, so it is
not in $\Wfin$. It is also the colimit of the sequence of finite coproducts $\coprod_{n<m}j$,
which lie in $\Wfin$.

\emph{Loc (9).} First, $\Wfin$ is closed under $\aleph_1$-filtered colimits. Let
$u=\colim_iu_i$ over an $\aleph_1$-filtered category, with $u_i\in\Wfin$. The first two
conditions of (fin) pass to the colimit as in Section~\ref{subsec:wtor}. If $u$ had infinitely
many bad components, choose countably many of them, and in each a based homotopy class
witnessing the failure, in some degree $n\ge2$: a non-zero class of the source mapped to zero,
or a class of the target outside the image. These countably many classes, with the null
homotopies of the first kind, come from one stage $i$. There the chosen components are still
distinct, a class of the first kind is still non-zero and still mapped to zero, and a class of
the second kind is still outside the image, since otherwise it would be in the image in the
colimit. So $u_i$ would have infinitely many bad components, which is impossible. Secondly, every
countable sub-functor of $u\in\Wfin$ lies in a countable sub-functor $u'$ that belongs to
$\Wfin$: run the construction of Section~\ref{subsec:wtor} with the requests of (fin) (the
fundamental groupoid and integral homology), with paths joining any two objects of the
sub-categories that lie in one component of $A$ or of $B$, and, on every component on which $u$
is a weak equivalence, with all requests on $\pi_n$, $n\ge2$. Then $\pi_0(u')$ is bijective,
each component of $B$ contains at most one component of $u'$, and $u'$ is a weak equivalence on
every component of $u'$ lying in a good component of $u$; so $u'$ has at most as many bad
components as $u$, and $u'\in\Wfin$. As for $\Wtor$, this is the form of Loc~(9) of
Section~\ref{subsec:axioms}. This completes the proof of Theorem~A.

\subsection{Countably many bad components: the class \texorpdfstring{$\Wctbl$}{W ctbl}}\label{subsec:wctbl}

Let $\Wctbl$ be the class of functors $u\colon A\to B$ such that $\Pione(Bu)$ is an
equivalence, $H_*(Bu;\Z)$ is an isomorphism, and $Bu$ restricts to a weak homotopy equivalence
on all but countably many connected components of $BA$; bad components are defined as in
Section~\ref{subsec:wfin}.

\begin{proposition}\label{prop:ctbl}
$\Wctbl$ satisfies Loc~\textup{(1)}, \textup{(2)}, \textup{(3)}; it is a strongly saturated weak
fundamental localizer that satisfies Loc~\textup{($6'$a)} and Loc~\textup{(8~bis)}, and it does
not satisfy Loc~\textup{(3~bis)}. So Loc~\textup{(8)} cannot be replaced by
Loc~\textup{(8~bis)} in Theorem~B\textup{(ii)}, nor in Corollaire~7.7.2.
\end{proposition}

\begin{proof}
Loc~(1), (2), (3) and Loc~($6'$a) are proved word for word as for $\Wfin$ in
Section~\ref{subsec:wfin}, with ``countably many'' in place of ``finitely many'': a union of two
countable sets is countable, an aspherical category is simply connected and acyclic, and in
the proof of Loc~($6'$a) the set $\bar S$ is now countable. Strong saturation follows by
Theorem~B(i). For Loc~(8~bis) let $u=\colim_nu_n$ over $\N$, with $u_n\colon A_n\to B_n$ in
$\Wctbl$; the first two conditions pass to the colimit as in Section~\ref{subsec:wtor}. Let $c$
be a component of $BB$, where $B=\colim_nB_n$, into which no bad component of any $u_n$ maps.
The objects of $B_n$ lying over $c$ form a union $B_n^c$ of components of $B_n$, the
corresponding components of $A_n$ form $A_n^c$, and $u_n$ restricts to a weak equivalence
$A_n^c\to B_n^c$, a disjoint union of weak equivalences. The parts of $A$ and $B$ over $c$ are
the colimits of the $A_n^c$ and of the $B_n^c$, and a sequential colimit of weak equivalences
is a weak equivalence, since homotopy groups commute with such colimits; so $c$ is not bad.
Hence every bad component of $u$ receives a bad component of some $u_n$, and $u$ has countably
many bad components. Finally, over the discrete category on an uncountable set $\Lambda$, the
coproduct of $\Lambda$ copies of the functor $j\colon Q\to P$ of Section~\ref{subsec:wfin} has
the comparison $j\in\Wfin\subseteq\Wctbl$ over every point and uncountably many bad
components, so Loc~(3~bis) fails. By Theorem~B(ii), $\Wctbl$ does not satisfy Loc~(8);
directly, that coproduct is the filtered colimit of the coproducts over the countable subsets
of $\Lambda$.
\end{proof}

\section{Gluing and filtered colimits give Loc~(3~bis)}\label{sec:positive}

This section proves Theorem~B. Throughout, $W$ is a weak fundamental localizer; by
\cite[Th.~6.1.18]{Cis06}, $\Winf\subseteq W$. For a simplicial set $X$ let $\Delta/X$ be its
category of simplices, and put
\[
V=\{f\colon X\to Y\ \text{in}\ \sSet:\ \Delta/f\in W\}.
\]
We use three facts from \cite{Cis06}. (F1) A map $f$ of simplicial sets is a weak homotopy
equivalence if and only if $\Delta/f\in\Winf$: by \cite[Th.~4.2.15]{Cis06}, $W'\mapsto
\{f:\Delta/f\in W'\}$ is the inverse of the bijection between test localizers of simplicial sets
and fundamental localizers modelable by $\Delta$, and the minimal fundamental localizer $\Winf$
corresponds to the minimal test localizer, the weak homotopy equivalences
\cite[Cor.~2.1.21 and proof of Cor.~4.2.19]{Cis06}. (F2) For every small category $A$ the
natural functor $\Delta/NA\to A$, sending a simplex to its last vertex, lies in $\Winf$
\cite[Th.~4.1.26 and Ex.~4.1.29]{Cis06}. (F3) By (F2), two out of three and
$\Winf\subseteq W$, a functor $u$ lies in $W$ if and only if $Nu\in V$.

By (F1), $V$ contains every weak homotopy equivalence; it satisfies two out of three because $W$
does, and it contains $i$ and $r$ whenever $ri=1$ and $ir\in V$, because $W$ is weakly
saturated.

\begin{lemma}\label{lem:pushout}
If $W$ satisfies Loc~\textup{($6'$a)}, then $V$ is stable under pushout along monomorphisms: if
$k\colon X\to Y$ is a monomorphism in $V$ and $g\colon X\to Z$ is any map, the map from $Z$ to
$T=Y\amalg_XZ$ lies in $V$.
\end{lemma}

\begin{proof}
Let $M=(X\times\Delta^1)\amalg_{X\times\{1\}}Z$ be the mapping cylinder of $g$, with the
monomorphisms $\iota_0\colon X\to M$ (at $X\times\{0\}$) and $\iota_Z\colon Z\to M$, and the
retraction $\rho\colon M\to Z$, so that $\rho\iota_0=g$ and $\rho\iota_Z=\id_Z$. The realization
of $M$ deformation retracts onto that of $Z$, so $\rho$ and $\iota_Z$ are weak homotopy
equivalences. Put $T'=Y\amalg_XM$, formed along the monomorphisms $k$ and $\iota_0$. Then $Y$ and
$M$ are simplicial subsets of $T'$ whose union is $T'$ and whose intersection is $X$.

For a simplicial subset $U\subseteq T'$ the category $\Delta/U$ is the full subcategory of
$\Delta/T'$ of the simplices lying in $U$, and it is a sieve, since the faces and degeneracies of
a simplex of $U$ lie in $U$. So $\Delta/Y$ and $\Delta/M$ form an MV square in $\Delta/T'$ with
intersection $\Delta/X$, and the inclusion of the intersection into $\Delta/Y$ is $\Delta/k\in W$.
Loc~($6'$a) gives $\Delta/M\to\Delta/T'$ in $W$.

The map $T'\to T$ induced by $\id_Y$ and $\rho$ is the pushout of $\rho$ along the monomorphism
$M\to T'$, since $T'\amalg_MZ=Y\amalg_XZ$. In the cofibration category of simplicial sets, with
the monomorphisms as cofibrations and the weak homotopy equivalences, every object is cofibrant,
and the pushout of a weak equivalence along a cofibration is a weak equivalence
\cite[Lemma~1.4.2]{RB09}; so $T'\to T$ is a weak homotopy equivalence. The map $Z\to T$ is the
composite $Z\xrightarrow{\iota_Z}M\to T'\to T$, and its image under $\Delta/-$ is the composite of
$\Delta/\iota_Z\in\Winf$, $\Delta/M\to\Delta/T'\in W$ and $\Delta/T'\to\Delta/T\in\Winf$ (by (F1)).
So $Z\to T$ lies in $V$.
\end{proof}

\begin{proof}[Proof of Theorem~B(i)]
Equip $\sSet$ with the monomorphisms as cofibrations and $V$ as weak equivalences. The axioms
CF1--CF4 of a precofibration category \cite[Def.~1.1.1]{RB09} hold: the initial object is
cofibrant, monomorphisms are closed under composition, isomorphisms lie in $V$, and isomorphisms
are trivial cofibrations (CF1); two out of three (CF2); pushouts of monomorphisms are
monomorphisms, and those of monomorphisms in $V$ lie in $V$ by Lemma~\ref{lem:pushout} (CF3);
every map $g$ factors as $\rho\iota_0$ through its mapping cylinder, with $\iota_0$ a
monomorphism and $\rho$ a weak homotopy equivalence, hence in $V$ (CF4). $V$ is weakly saturated
in the sense of \cite[Def.~7.2.6]{RB09}, as noted above. By \cite[Th.~7.2.7]{RB09}, $V$ is
saturated: a map lies in $V$ as soon as it becomes invertible in $\sSet[V^{-1}]$.

Let $u$ be a functor whose image in $\Cat[W^{-1}]$ is invertible. By (F3) the nerve sends $W$
into $V$, so it induces a functor $\Cat[W^{-1}]\to\sSet[V^{-1}]$, and the image of $Nu$ in
$\sSet[V^{-1}]$ is invertible. So $Nu\in V$, and $u\in W$ by (F3).
\end{proof}

\begin{proof}[Proof of Theorem~B(ii), sufficiency]
Assume Loc~($6'$a) and Loc~(8). We show that $V$ is a localizer of simplicial sets in the sense of
\cite[D\'ef.~1.4.1]{Cis06}: (L1) two out of three; (L2) every trivial fibration lies in $V$;
(L3) the monomorphisms in $V$ are stable under pushouts and transfinite compositions. (L1) and
(L2) hold since $V$ contains the weak homotopy equivalences. Stability under pushouts is
Lemma~\ref{lem:pushout}. Let $X_0\to X_1\to\cdots\to X_\beta\to\cdots$ ($\beta<\lambda$) be a
transfinite sequence of monomorphisms in $V$, continuous at limit ordinals (the form of \cite[Rem.~1.1.3]{Cis06}); we show by induction
that $X_0\to X_\beta$ lies in $V$ for all $\beta\le\lambda$. Successor steps follow from two out of
three. At a limit ordinal $\beta$, the arrow $X_0\to X_\beta$ is the colimit, over the filtered
set of $\beta'<\beta$, of the arrows $X_0\to X_{\beta'}$; the functor $X\mapsto\Delta/X$ commutes
with filtered colimits, because simplices, the maps between them and their relations are finite
data; and Loc~(8) gives $\Delta/X_0\to\Delta/X_\beta$ in $W$.

So $V$ is a localizer of simplicial sets; it is regular \cite[Prop.~3.4.25]{Cis06}, and since the
maps $\Delta^n\to\Delta^0$ are weak equivalences it is a test localizer
\cite[D\'ef.~4.2.11]{Cis06}. By \cite[Th.~4.2.15 and Rem.~4.2.16]{Cis06}, applied with the
inclusion of $\Delta$ into $\Cat$, whose associated functor is the nerve \cite[Ex.~4.1.29]{Cis06},
the class $N^{-1}V$ of functors $u$ with $Nu\in V$ is a fundamental localizer. By (F3),
$N^{-1}V=W$.
\end{proof}

\begin{proposition}\label{prop:necessity}
Every fundamental localizer satisfies Loc~\textup{($6'$a)} and Loc~\textup{(8)}.
\end{proposition}

\begin{proof}
Let $W$ be a fundamental localizer. Loc~(8) is \cite[Cor.~4.2.22]{Cis06}. The category $\Delta$ is
a $W$-test category \cite[Ex.~4.1.21]{Cis06}, so $V=\{f:\Delta/f\in W\}$ is a regular localizer of
simplicial sets \cite[Prop.~4.2.9]{Cis06}, and $W=N^{-1}V$ \cite[Th.~4.1.26 and
Ex.~4.1.29]{Cis06}. For an MV square $X=X_1\cup X_2$, Lemma~\ref{lem:mv} says that
$NX=NX_1\amalg_{NX_0}NX_2$. If $X_0\to X_1$ lies in $W$, then $NX_0\to NX_1$ is a monomorphism in
$V$, so is its pushout $NX_2\to NX$ by axiom (L3) of \cite[D\'ef.~1.4.1]{Cis06}, and $X_2\to X$
lies in $N^{-1}V=W$.
\end{proof}

Proposition~\ref{prop:necessity} is the necessity in Theorem~B(ii), which is now proved.

\section{The passage of \emph{Les D\'erivateurs}}\label{sec:manuscript}

We cite the edition of Chapter~XVI, \emph{Les localiseurs fondamentaux}, deciphered and
transcribed by M.~K\"unzer and edited by M.~K\"unzer, J.~Malgoire and G.~Maltsiniotis, in its
version of 18~March 2019 \cite{Der}; page numbers are those printed in the edition, and the
manuscript pages are given in brackets in the edition. Loc~(1), (2), (3) are stated on p.~1
(ms.~pp.~1--2), together with footnote~1 on the exchange of Loc~(3) and Loc~(3~bis);
D\'efinition~1.1 and Loc~(3~bis) are on p.~2 (ms.~pp.~2--3), the MV axiom Loc~($6'$) on p.~13
(ms.~p.~$17'$), Loc~(8) and Loc~(9) in \S5, on pp.~40--41.

In \S7.7 Grothendieck shows that a closed model category given by a special Quillen triple
defines a prederivator whose localizer on $\Cat$ is a localiseur fondamental satisfying
Loc~(1~ter, 3~bis), and a derivator when the triple is very special (Proposition~7.7.1). He continues: if a localiseur fondamental $W$ satisfies Loc~($6'$a,~8,~9),
then $\mathrm{HOT}_W$ is a derivator, and ``On a vu de plus dans 7.6 (page 111) que l'on a alors
$W=W_{\mathrm{HOT}_W}$'', with footnote~55: ``Mais c'\'etait sous l'hypoth\`ese Loc (3 bis)~!
Gare au cercle vicieux.'' Then, on p.~55, lines~16--17 (counting every line of the page below the running head,
including the editors' manuscript-page markers; page~57 of the edition's file for the chapter),
manuscript page~118:
\begin{quote}
\textbf{Corollaire 7.7.2~:} \emph{On a l'implication}
\[
\mathrm{Loc}\,(6'\ \mathrm{a},8,9)\ \Longrightarrow\ \mathrm{Loc}\,(1\ \mathrm{ter},3\ \mathrm{bis})\ {}^{(56)}.
\]
\end{quote}
and at the foot of the same page, footnote~56: ``Loc (3 bis) pas prouv\'e, probablement faux.''

Corollary~C gives the status of each part. The implication is true (Theorem~B), and the
circularity named in footnote~55 is avoided because the proof of Theorem~B never identifies $W$
with the localizer of a derivator before Loc~(3~bis) is established. So footnote~56 is right in
its first half and wrong in its second. The argument of \S7.7 does not prove Loc~(3~bis); but
under the hypotheses of the corollary Loc~(3~bis) is not false: it holds, and so does
Loc~(1~ter) (Theorem~B(i)), and Loc~(9) plays no part. The doubt is borne out only for weaker
hypotheses: Loc~(3~bis) fails as soon as Loc~($6'$a) is dropped, even with Loc~(1~ter), (8) and
(9) (the class $\Wtor$), and as soon as Loc~(8) is dropped, even with Loc~(1~ter), ($6'$a) and
(9) (the class $\Wfin$), or is weakened to Loc~(8~bis) (the class $\Wctbl$). This also settles the hesitation of footnote~57 of \S7.8, whether to add Loc~(3~bis) to
the definition of the \emph{localiseurs sp\'eciaux} (those satisfying Loc~($6'$a,~8,~9)): they
satisfy it automatically.

\begin{proof}[Proof of Corollary~C]
A class satisfying Loc~(1), (2), (3) is a weak fundamental localizer
(Section~\ref{subsec:axioms}). Part~(i) is Theorem~B. Part~(ii) is Theorem~A(c) and~(d). Part~(iii)
is Theorem~A(a).
\end{proof}

\section{What remains open}\label{sec:open}

\begin{enumerate}
\item \emph{Coproducts.} Loc~(8) cannot be replaced by Loc~(8~bis) in Theorem~B
(Proposition~\ref{prop:ctbl}). Does a weak fundamental localizer satisfying Loc~($6'$a) and
stability under arbitrary coproducts satisfy Loc~(3~bis)? $\Wfin$ and $\Wctbl$ are not stable
under arbitrary coproducts.
\item \emph{Strong saturation without gluing.} Does a weak fundamental localizer satisfying
Loc~(8) and Loc~(9) satisfy Loc~(1~ter)? It satisfies Loc~(1~bis). It is closed under retracts
by Lemma~\ref{lem:retracts}. It satisfies Loc~(1)(c): if $f\colon X\to Y$ and $g\colon Y\to X$
have $gf$ and $fg$ in $W$, the canonical maps from $X$ and from $Y$ to the colimits $X_\infty$
of $X\xrightarrow{gf}X\xrightarrow{gf}\cdots$ and $Y_\infty$ of
$Y\xrightarrow{fg}Y\xrightarrow{fg}\cdots$ lie in $W$ by Loc~(8~bis), being the colimits of the
arrows $(gf)^n$ and $(fg)^n$; the map $X_\infty\to Y_\infty$ induced by $f$ is an isomorphism,
with inverse induced by $g$; so $f\in W$ by two out of three, and likewise $g$. The class
$\Wtor$, which fails Loc~($6'$a), is strongly saturated; we have no example either way.
\item \emph{The remaining entries of Table~\ref{tab:classes}}: Loc~(9) for $\Wcirc$ and for
$\Wctbl$.
\end{enumerate}

\section{Note on method}\label{sec:method}

This paper is AI-produced. No human has proved, written or checked any part of it.

\paragraph{What is whose.} The axioms and Corollaire~7.7.2 with its footnotes are Grothendieck's;
the weak notion is Maltsiniotis's; the results of \cite{Cis06} named in
Section~\ref{subsec:credit}, on which Theorem~B and Proposition~\ref{prop:necessity} rest, are
Cisinski's, and the saturation theorem used for Theorem~B(i) is Rădulescu-Banu's. The new
mathematics was found by GPT-6 Astra (OpenAI, ChatGPT Pro mode) in four runs, between 11 and
18~November 2026: the class $\Wpi$ with its witness on 11~November, first, and $\Wcirc$ with its
witness in another run later that day; $\Wtor$ and the sufficiency in Theorem~B(ii) in a third run
(by 16~November); and $\Wfin$ with Theorem~B(i) in a fourth (18~November). The text was written by Claude Opus~5.5 at effort max, who assembled
the four runs into one argument and renamed their classes. The following are the writer's: the
formulation through Lemmas~\ref{lem:absolute}--\ref{lem:mv}; Lemma~\ref{lem:retracts}; the
failure of Loc~($6'$a) for $\Wpi$ and $\Wcirc$; the mapping-cylinder proof of
Lemma~\ref{lem:pushout}, which replaces a comparison between a double collage and a category of
simplices; and Proposition~\ref{prop:necessity}, deduced from Cisinski's results.
Proposition~\ref{prop:ctbl} (the class $\Wctbl$), the failure of Loc~(8~bis) for $\Wcirc$
and the failure of Loc~(9) for $\Wpi$ were found by Claude Fable~5.2 at effort max while
checking this paper.

\paragraph{Checks.} Before this text was written, each of the four separating classes and its
witness was followed step by step by a Claude model: $\Wpi$ and $\Wtor$ by Claude Opus~5.5 at
effort high (24~November 2026), $\Wtor$ also by Claude Opus~5.5 at effort max (18~November),
and $\Wcirc$ and $\Wfin$ by Claude Opus~5.5 at effort max (19~November), who also proved the
strong saturation of $\Wcirc$ again independently; the homological step of the witness for
$\Wfin$ was derived again on 24~November (Claude Opus~5.5 at effort high). The gluing step of
Theorem~B(i) was followed, and the sufficiency in Theorem~B(ii) audited, by Claude Opus~5.5 at
effort max (19~November), without opening the cited theorems of Cisinski and
Rădulescu-Banu. The writer opened every cited statement of \cite{Cis06}, \cite{RB09}, \cite{Der} and
\cite{DerXVIII} and derived every proof in this text again while writing it. The runs verified each
finite witness with a program (\texttt{a\_absolute\_relative.py},
\texttt{verify\_counterexamples.py}, \texttt{FGH\_literal\_axioms.py},
\texttt{A5\_finite\_component\_counterexample.py}). All these checks are by models of a family
other than the one that found the results; the arguments that are the writer's own were checked line by
line on 25~November 2026 by Claude Fable~5.2 at effort max, which also checked
Section~\ref{sec:positive} against the cited statements of \cite{Cis06} and \cite{RB09},
opened \cite{Bro06}, \cite{Hat02}, \cite{Qui73} and \cite{Mal05} at the places cited, and
recounted the finite witnesses with a program of its own. That checker is a model of the same
family as the writer. Its three results named above were checked line by line by the
writer. Nothing is formalized. The record of every check, with model, effort and
date, is in the file \texttt{CHECK\_EVENTS.tsv} that accompanies this paper in its workbench,
reached from \url{https://github.com/KokunoYumeto/mathematics-workbenches};
\texttt{PROVENANCE.md} lists every contributor and says where each citation was confirmed.

\begin{thebibliography}{99}
\bibitem{Bro06} R.~Brown, \emph{Topology and groupoids}, third edition, BookSurge, 2006.
\bibitem{Cis06} D.-C.~Cisinski, \emph{Les pr\'efaisceaux comme mod\`eles des types
 d'homotopie}, Ast\'erisque 308, Soc.\ Math.\ France, 2006.
\bibitem{Der} A.~Grothendieck, \emph{Les D\'erivateurs}, manuscript, 1990--1991; Chapitre~XVI,
 \emph{Les localiseurs fondamentaux}, deciphered and transcribed by M.~K\"unzer, edited by
 M.~K\"unzer, J.~Malgoire and G.~Maltsiniotis, version of 18~March 2019,
 \url{https://webusers.imj-prg.fr/~georges.maltsiniotis/groth/Derivateurs.html}.
\bibitem{DerXVIII} A.~Grothendieck, \emph{Les D\'erivateurs}, Chapitre~XVIII, \emph{Cat\'egories
 et ensembles accessibles}, same edition, version of 17~December 2025, pp.~72 and 104--105.
\bibitem{Hat02} A.~Hatcher, \emph{Algebraic topology}, Cambridge University Press, 2002.
\bibitem{Mal05} G.~Maltsiniotis, \emph{La th\'eorie de l'homotopie de Grothendieck},
 Ast\'erisque 301, Soc.\ Math.\ France, 2005.
\bibitem{Qui73} D.~Quillen, Higher algebraic $K$-theory~I, in \emph{Algebraic $K$-theory~I},
 Lecture Notes in Math.\ 341, Springer, 1973, 85--147.
\bibitem{RB09} A.~R\u{a}dulescu-Banu, Cofibrations in homotopy theory, arXiv:math/0610009v4,
 2009.
\bibitem{Tho79} R.~W.~Thomason, Homotopy colimits in the category of small categories,
 Math.\ Proc.\ Cambridge Philos.\ Soc.\ 85 (1979), 91--109.
\end{thebibliography}
\end{document}
